%% toc: P15 158 158.1,158.2,158.5,158.8
\chapter{Endpoint Limits, Stieltjes Laws and Trace Sigma-Fields}
\label{sec:p15-w02c-p15-003-stieltjes-endpoints}

The construction of the measure \(\mu_F\) of a Stieltjes measure function
in \cref{thm:p15-stieltjes-construction} rests on three facts about the
function and its measure. The endpoint values \(F(-\infty)\) and
\(F(\infty)\), defined as an infimum and a supremum, are the limits of
\(F\) at the two ends of the line, also when they are infinite. The masses
of the bounded truncations of an unbounded interval increase to the mass
assigned to the whole interval. And every bounded interval, whichever of
its endpoints it contains, has finite measure. The quantile theorem
\cref{thm:p15-cdf-existence} adds two conclusions about laws: \(\mu_F\)
is the law of the quantile variable \(Q_F\), and it is the only Borel
probability law whose cumulative distribution function is \(F\). Those
conclusions need the fact that a Borel probability measure on the line is
the law of a real random variable and is determined by its values on the
closed rays. A third section concerns the density construction of
\cref{thm:p15-countably-generated-radon-nikodym}: its sigma-finite case
applies the finite construction on each block of a measurable partition, and
that needs the restricted sigma-field on a block to be generated by the
restrictions of the original generating sets, and it needs integrals
against the normalised measure on a block to be converted into integrals
against the original measure on the whole space before the local densities
are glued. This chapter states and
proves all of these facts.

\emph{Inputs.} A Stieltjes measure function \(F\) in the sense of
\cref{def:p15-stieltjes-intervals}, with its measure \(\mu_F\); for the
second section, a Borel probability measure on \(\mathbb R\), or a
function \(F\) satisfying the hypotheses of \cref{def:p15-quantile-map}.
\emph{Conclusions.} The two endpoint limits; the truncation identity for
interval masses with infinite endpoints or infinite endpoint values;
finiteness of \(\mu_F\) on bounded Borel sets, with the exact measure of
each kind of bounded interval; the identification of a Borel probability
measure as the law of its identity variable; uniqueness of a Borel
probability measure given its values on closed rays; and the resulting
identification of \(\mu_F\) with the law of \(Q_F\) and with every Borel
probability law whose cumulative distribution function is \(F\).

\emph{Source.} Durrett's author version of \emph{Probability: Theory and
Examples}, dated 11 January 2019, states the Stieltjes existence and
uniqueness theorem as Theorem~1.1.4 \cite[thm.~1.1.4, p.~3]{DurrettPTEAuthor2019}.
At the start of its proof he records that the limits of \(F\) at
\(\infty\) and at \(-\infty\) exist, and that the interval assignment
is meaningful for every allowed pair of endpoints because those limits are
not \(-\infty\) and \(\infty\) respectively
\cite[p.~5]{DurrettPTEAuthor2019}. The existence of those limits as the
supremum and infimum is proved below. His Theorem~1.2.2 constructs a real
random variable with a prescribed distribution function by the generalized
inverse \cite[thm.~1.2.2, p.~11]{DurrettPTEAuthor2019}, and the remark
following it derives, from the uniqueness clause of Theorem~1.1.4, that two
real random variables have the same law exactly when they have the same
distribution function \cite[pp.~11--12]{DurrettPTEAuthor2019}. The second
section proves the corresponding statements for Borel probability measures
that are not presented as laws of given random variables.

\section{Endpoint limits, truncations and bounded intervals}
\label{sec:p15-w02c-p15-003-endpoint-limits}

\begin{hbdefinition}{def:p15-w02c-p15-003-extended-limits}{Extended limits at the ends, left limits and truncations}
Let \(F\) be a Stieltjes measure function. For
\(L\in(-\infty,\infty]\), say that \(F(x)\to L\) as \(x\to\infty\)
when either \(L\) is finite and for every \(\varepsilon>0\) there is a
real \(x_0\) with \(|F(x)-L|<\varepsilon\) for all \(x\geq x_0\), or
\(L=\infty\) and for every real \(M\) there is a real \(x_0\) with
\(F(x)>M\) for all \(x\geq x_0\). For \(L\in[-\infty,\infty)\),
the statement \(F(x)\to L\) as \(x\to-\infty\) is defined in the same
way with \(x\leq x_0\) in place of \(x\geq x_0\), and, when
\(L=-\infty\), with \(F(x)<M\) in place of \(F(x)>M\).

For real \(c\), the \emph{left limit} is
\[
 F(c-)\defeq\sup_{t<c}F(t).
\]
It is a real number, since \(F(c-1)\leq F(c-)\leq F(c)\). This extends
the notation of \cref{def:p15-cumulative-distribution} from bounded
functions to Stieltjes measure functions.

For \(-\infty\leq a<b\leq\infty\) and integers \(n\geq1\), the
\emph{truncations} of \(I(a,b)\) are
\[
 K_n\defeq I(a,b)\cap(-n,n],\qquad
 l_n\defeq\max(a,-n),\qquad r_n\defeq\min(b,n).
\]
The numbers \(l_n\) and \(r_n\) are real, because \(a<\infty\) and
\(b>-\infty\).
\index{Stieltjes measure!endpoint limits}\index{interval!truncation}
\end{hbdefinition}

\begin{hblemma}{thm:p15-w02c-p15-003-endpoint-limits}{proved=annex:p15-w02c-p15-003-endpoint-limits}{Endpoint limits and the truncation identity}
Let \(F\) be a Stieltjes measure function.
\begin{enumerate}
\item The endpoint values are the limits at the two ends:
      \(F(x)\to F(\infty)\) as \(x\to\infty\), and
      \(F(x)\to F(-\infty)\) as \(x\to-\infty\), in the sense of
      \cref{def:p15-w02c-p15-003-extended-limits}, including the cases
      \(F(\infty)=\infty\) and \(F(-\infty)=-\infty\). The sequence
      \(F(n)\) is nondecreasing with supremum \(F(\infty)\), and the
      sequence \(F(-n)\) is nonincreasing with infimum \(F(-\infty)\).
\item Let \(-\infty\leq a<b\leq\infty\). There is a least integer
      \(n_0\geq1\) with \(l_{n_0}<r_{n_0}\). The truncation \(K_n\) is
      empty for \(n<n_0\) and equals \((l_n,r_n]\) for
      \(n\geq n_0\). The masses \(m_F(K_n)\) are real, nondecreasing in
      \(n\), and
      \[
       \sup_{n\geq1}m_F(K_n)=m_F(I(a,b))=F(b)-F(a),
      \]
      where the right side is evaluated by the endpoint convention of
      \cref{def:p15-stieltjes-intervals} and may be \(\infty\).
\end{enumerate}
\end{hblemma}

Part~(2) is the truncation identity used in the proof of
\cref{thm:p15-stieltjes-construction} to pass from bounded targets to
targets with an infinite endpoint. Its proof distinguishes three cases:
both endpoint values finite, \(F(b)=\infty\), and \(F(a)=-\infty\). It
never subtracts one infinite value from another.

\begin{hbcorollary}{thm:p15-w02c-p15-003-bounded-finite}{proved=annex:p15-w02c-p15-003-bounded-finite}{Finite mass on bounded sets and the measure of each bounded interval}
Let \(F\) be a Stieltjes measure function and \(\mu_F\) its measure.
If \(B\) is a Borel set with \(B\subseteq[-r,r]\) for some real
\(r\), then
\[
 \mu_F(B)\leq F(n)-F(-n)<\infty
 \qquad\text{for every integer }n>r .
\]
Every bounded interval is a bounded Borel set, so its measure is finite.
For real \(c<d\), the measures are
\[
 \begin{aligned}
 \mu_F(\{c\})&=F(c)-F(c-), &\qquad
 \mu_F((c,d])&=F(d)-F(c),\\
 \mu_F((c,d))&=F(d-)-F(c), &\qquad
 \mu_F([c,d])&=F(d)-F(c-),\\
 \mu_F([c,d))&=F(d-)-F(c-). &&
 \end{aligned}
\]
\end{hbcorollary}

The bound requires boundedness of the set, not of the function: for
\(F(x)=x\), every truncation \((-n,n]\) has finite measure \(2n\),
while \(\mathbb R\) has measure \(\infty\). Boundedness is sufficient
for finiteness but not necessary, as the next example shows.

\begin{hbexample}{def:p15-w02c-p15-003-exponential}{An infinite upper endpoint value}
Let \(F(x)=e^x\). It is real valued, increasing and continuous, so it is a
Stieltjes measure function. Its infimum \(F(-\infty)=0\) is not
attained, and its supremum is \(F(\infty)=\infty\).

For the target \(I(0,\infty)=(0,\infty)\), the truncations are
\(K_n=(0,n]\) with masses \(e^n-1\). They increase without bound, and the
assigned mass is \(F(\infty)-F(0)=\infty-1=\infty\), as part~(2) of
\cref{thm:p15-w02c-p15-003-endpoint-limits} requires. For
\(I(-\infty,\infty)=\mathbb R\), the truncations \((-n,n]\) have masses
\(e^n-e^{-n}\), which also increase to
\(\mu_F(\mathbb R)=F(\infty)-F(-\infty)=\infty-0=\infty\).

Because \(F\) is continuous, \(F(c-)=F(c)\) for every \(c\), and
\cref{thm:p15-w02c-p15-003-bounded-finite} gives every singleton measure
zero and every bounded interval with endpoints \(c<d\) the measure
\(e^d-e^c\), whichever endpoints it contains. The unbounded ray
\((-\infty,0]\) has the finite measure \(F(0)-F(-\infty)=1\).
\end{hbexample}

\begin{hbwarning}{def:p15-w02c-p15-003-monotone-premise}{The endpoint limits use monotonicity}
The supremum of a real function need not be its limit at infinity. The
function \(\sin x\) has supremum one, but it has no limit as
\(x\to\infty\): it equals one at \(x=\pi/2+2k\pi\) and minus one at
\(x=3\pi/2+2k\pi\) for every integer \(k\geq0\). Part~(1) of
\cref{thm:p15-w02c-p15-003-endpoint-limits} uses that \(F\) is
nondecreasing at every step, and the truncation identity depends on it.
Finiteness on bounded intervals also does not pass to unbounded ones:
for \(F(x)=x\), every bounded interval has finite length and
\(\mu_F((0,\infty))=\infty\).
\end{hbwarning}

\section{Laws determined by closed rays}
\label{sec:p15-w02c-p15-003-ray-laws}

A cumulative distribution function is defined in
\cref{def:p15-cumulative-distribution} for a real random variable, and
\cref{thm:p15-cdf-laws} compares the laws of two real random variables.
A Borel probability measure on the line that is not presented as the law of
a given random variable is brought under those definitions by its identity
variable, and is determined by its values on the closed rays.

\begin{hbtheorem}{thm:p15-w02c-p15-003-ray-laws}{proved=annex:p15-w02c-p15-003-ray-laws}{A Borel probability is a law and is determined by its rays}
\par\nobreak
\begin{enumerate}
\item Let \(\nu\) be a probability measure on
      \((\mathbb R,\mathcal B(\mathbb R))\). On the probability space
      \((\mathbb R,\mathcal B(\mathbb R),\nu)\), the identity map
      \(\iota(x)=x\) is a real random variable, its law is \(\nu\), and
      its cumulative distribution function is
      \(F_\iota(x)=\nu((-\infty,x])\).
\item If \(\nu_1,\nu_2\) are probability measures on
      \((\mathbb R,\mathcal B(\mathbb R))\) and
      \(\nu_1((-\infty,x])=\nu_2((-\infty,x])\) for every real \(x\),
      then \(\nu_1=\nu_2\).
\item Let \(F\) satisfy the hypotheses of \cref{def:p15-quantile-map}.
      Then \(F(-\infty)=0\), \(F(\infty)=1\), \(\mu_F\) is a
      probability measure with \(\mu_F((-\infty,x])=F(x)\), and
      \(\mu_F\) is the law of the identity variable on
      \((\mathbb R,\mathcal B(\mathbb R),\mu_F)\), whose cumulative
      distribution function is \(F\). The law of \(Q_F\) on the uniform
      probability space is \(\mu_F\). A probability measure \(\nu\) on
      \((\mathbb R,\mathcal B(\mathbb R))\) satisfies
      \(\nu((-\infty,x])=F(x)\) for every real \(x\) if and only if
      \(\nu=\mu_F\). A real random variable \(X\) on any probability
      space satisfies \(F_X=F\) if and only if its law is \(\mu_F\).
\end{enumerate}
\end{hbtheorem}

Part~(3) is the step between the ray identity
\(\mathbb U(Q_F\leq x)=F(x)\) and the conclusions of
\cref{thm:p15-cdf-existence} that the law of \(Q_F\) is \(\mu_F\) and that
\(\mu_F\) is the unique Borel probability law with cumulative distribution
function \(F\). Part~(2) is proved directly from the closed rays as a
generating pi-system, and also, as in Durrett's remark, from the comparison
of distribution functions in \cref{thm:p15-cdf-laws} applied to the identity
variables of part~(1).

\begin{hbexample}{def:p15-w02c-p15-003-two-point}{A two-point law from its cumulative function}
Let \(F(x)=0\) for \(x<0\), \(F(x)=1/2\) for \(0\leq x<1\), and
\(F(x)=1\) for \(x\geq1\). It satisfies the hypotheses of
\cref{def:p15-quantile-map}. For real \(c\), let \(\delta_c\) be the
Borel measure with \(\delta_c(B)=1\) when \(c\in B\) and zero otherwise,
and let \(\nu=\tfrac12\delta_0+\tfrac12\delta_1\). This \(\nu\) is a
Borel probability measure, and \(\nu((-\infty,x])\) is \(0\) for
\(x<0\), \(1/2\) for \(0\leq x<1\) and \(1\) for \(x\geq1\). By part~(3)
of \cref{thm:p15-w02c-p15-003-ray-laws}, \(\nu=\mu_F\). In agreement,
\cref{thm:p15-w02c-p15-003-bounded-finite} gives
\(\mu_F(\{0\})=F(0)-F(0-)=1/2\) and \(\mu_F(\{1\})=1-1/2=1/2\).

The quantile is \(Q_F(u)=0\) for \(0<u\leq1/2\) and \(Q_F(u)=1\) for
\(1/2<u<1\), since \(\{y:F(y)<u\}\) is \((-\infty,0)\) in the first
case and \((-\infty,1)\) in the second. Its law gives each of \(0\) and
\(1\) uniform probability \(1/2\), which is \(\mu_F\) again.
\end{hbexample}

\begin{hbwarning}{def:p15-w02c-p15-003-ray-limits}{What the ray values determine}
Part~(2) of \cref{thm:p15-w02c-p15-003-ray-laws} is stated for probability
measures and uses the values on a generating family of rays. Two infinite
measures can agree on every closed ray and still differ: let
\(\nu_1(B)\) be the number of integers in \(B\), with value
\(\infty\) when \(B\) contains infinitely many integers, and let
\(\nu_2=2\nu_1\). Both are Borel measures, since the number of integers
in a disjoint union is the sum of the numbers in its members. Every ray \((-\infty,x]\) contains infinitely
many integers, so both measures give it the value \(\infty\), yet
\(\nu_1(\{0\})=1\) and \(\nu_2(\{0\})=2\).

Agreement on too few rays also fails. The probability measures
\(\delta_{1/4}\) and \(\delta_{1/2}\) give the ray \((-\infty,k]\) the
same value for every integer \(k\), namely zero when \(k\leq0\) and one
when \(k\geq1\), but they differ on \((-\infty,1/4]\). The integer rays do
not generate the Borel sets. The rational rays do, by
\cref{thm:p15-borel-generators}, and they form a pi-system because
\((-\infty,p]\cap(-\infty,q]=(-\infty,\min(p,q)]\); so
\cref{thm:p15-measure-uniqueness} gives equality of two probability
measures that agree on every rational ray.
\end{hbwarning}

\emph{Relationships.} The lemma of the first section supplies the endpoint
behaviour that \cref{thm:p15-stieltjes-construction} uses for its covering
bound on unbounded targets, for its second partition condition, and for the
interval masses \(\mu_F(I(a,\infty))=F(\infty)-F(a)\) and
\(\mu_F(\mathbb R)=F(\infty)-F(-\infty)\). Its corollary is the finiteness
on bounded intervals asserted there. The theorem of the second section
connects that construction with \cref{thm:p15-cdf-existence} and with the
pushforward laws of \cref{def:p15-pushforward-law}.

\section{Traces of countably generated sigma-fields}
\label{sec:p15-w02c-p15-003-trace-sigma-fields}

The density construction of \cref{thm:p15-countably-generated-radon-nikodym}
is first carried out for a finite measure and a probability measure on a
sigma-field generated by a sequence of sets \(E_1,E_2,\ldots\). Its
sigma-finite case cuts \(\Omega\) into countably many measurable blocks of
finite mass and runs the finite construction on each block. That step treats
a block \(D\) as a measurable space of its own, with the sets of the form
\(A\cap D\), and it needs this space to be generated by the sequence
\(E_1\cap D,E_2\cap D,\ldots\). The present section proves that equality,
by an inclusion in each direction, together with the facts about restricted
measures and glued functions that the assembly also uses, the transfer of
integrals from a block to the whole space, and the assembly of the local
densities into one density on \(\Omega\).

\emph{Inputs.} A measurable space \((\Omega,\mathcal F)\), usually with
\(\mathcal F=\sigma(E_1,E_2,\ldots)\); a nonempty subset \(D\) of
\(\Omega\), for the measure statements a member of \(\mathcal F\); measures
\(\mu,\nu\) on \((\Omega,\mathcal F)\) with \(\mu\ll\nu\) in the sense of
\cref{def:p15-absolute-continuity}. \emph{Conclusions.} The restricted
collection on \(D\) is a sigma-field on \(D\); when \(D\in\mathcal F\) its
members are exactly the events contained in \(D\), and restricted measures
are measures; the restriction of a generated sigma-field is generated by the
restricted generators; functions glued from measurable pieces on a countable
measurable partition are measurable; a block of finite positive mass
carries exactly the data required by the finite case of the density
construction; integrals on such a block against its normalised measure are
\(\nu(D)^{-1}\) times the \(\nu\)-integrals over \(\Omega\) of the zero
extensions, with the same null sets; and local densities on the blocks glue
to a density of \(\mu\) relative to \(\nu\) on all of \(\Omega\),
finite \(\nu\)-almost everywhere.

\begin{hbdefinition}{def:p15-w02c-p15-003-trace-sigma-field}{Trace sigma-fields, sigma-fields generated on a subset, restricted measures}
Let \((\Omega,\mathcal F)\) be a measurable space and let \(D\) be a
nonempty subset of \(\Omega\). The \emph{trace} of \(\mathcal F\) on
\(D\) is
\[
              \mathcal F_D\defeq\{A\cap D:A\in\mathcal F\}.
\]
For a collection \(\mathcal C\subseteq2^D\), the sigma-field
\(\sigma_D(\mathcal C)\) is the collection \(\sigma(\mathcal C)\) of
\cref{def:p15-sigma-algebra} formed with \(D\) in place of
\(\Omega\): the intersection of all sigma-algebras on \(D\) that contain
\(\mathcal C\). Complements in a sigma-algebra on \(D\) are taken
relative to \(D\), so that the complement of \(B\subseteq D\) is
\(D\setminus B\). For sets \(C_1,C_2,\ldots\subseteq D\), write
\(\sigma_D(C_1,C_2,\ldots)\) for \(\sigma_D(\{C_j:j\geq1\})\), and
likewise \(\sigma(E_1,E_2,\ldots)\) for
\(\sigma(\{E_j:j\geq1\})\) on \(\Omega\).
When \(D\in\mathcal F\) and \(\mu\) is a measure on
\((\Omega,\mathcal F)\), the \emph{restriction} \(\mu|_D\) is the
function on \(\mathcal F_D\) given by \(\mu|_D(B)\defeq\mu(B)\); this
is meaningful because \(\mathcal F_D\subseteq\mathcal F\) in that case,
by part~(2) of \cref{thm:p15-w02c-p15-003-trace-generation}.
\index{sigma-algebra!trace}\index{measure!restriction}
\end{hbdefinition}

\begin{hbtheorem}{thm:p15-w02c-p15-003-trace-generation}{proved=annex:p15-w02c-p15-003-trace-generation}{Traces of generated sigma-fields}
Let \((\Omega,\mathcal F)\) be a measurable space and let \(D\) be a
nonempty subset of \(\Omega\).
\begin{enumerate}
\item The trace \(\mathcal F_D\) is a sigma-algebra on \(D\).
\item If \(D\in\mathcal F\), then
      \(\mathcal F_D=\{A\in\mathcal F:A\subseteq D\}\). For every measure
      \(\mu\) on \((\Omega,\mathcal F)\), the restriction \(\mu|_D\) is a
      measure on \((D,\mathcal F_D)\), and it is finite when
      \(\mu(D)<\infty\).
\item If \(\mathcal F=\sigma(\mathcal C)\) for a collection
      \(\mathcal C\subseteq2^\Omega\), then
      \[
           \mathcal F_D=\sigma_D(\{C\cap D:C\in\mathcal C\}).
      \]
      In particular, if \(\mathcal F=\sigma(E_1,E_2,\ldots)\), then
      \(\mathcal F_D=\sigma_D(E_1\cap D,E_2\cap D,\ldots)\).
\item Let \((D_i)_{i\in I}\) be a countable family of pairwise disjoint
      members of \(\mathcal F\) whose union is \(\Omega\), and let
      \((S,\mathcal S)\) be a measurable space. Suppose that for every
      \(i\) with \(D_i\neq\varnothing\) a map \(g_i:D_i\to S\) is
      measurable from \((D_i,\mathcal F_{D_i})\). Then the map
      \(g:\Omega\to S\) equal to \(g_i\) on each nonempty \(D_i\) is
      measurable from \((\Omega,\mathcal F)\).
\end{enumerate}
\end{hbtheorem}

\begin{hbcorollary}{thm:p15-w02c-p15-003-block-data}{proved=annex:p15-w02c-p15-003-block-data}{Blocks of finite positive mass}
Let \(\mathcal F=\sigma(E_1,E_2,\ldots)\) on \(\Omega\), let
\(\mu,\nu\) be measures on \((\Omega,\mathcal F)\) with \(\mu\ll\nu\),
and let \(D\in\mathcal F\) satisfy \(\mu(D)<\infty\) and
\(0<\nu(D)<\infty\). Then \(D\) is nonempty, and
\begin{enumerate}
\item \((D,\mathcal F_D)\) is a measurable space with
      \(\mathcal F_D=\sigma_D(E_1\cap D,E_2\cap D,\ldots)\);
\item \(\mu|_D\) is a finite measure on \((D,\mathcal F_D)\), and
      \(\lambda\defeq\nu|_D/\nu(D)\), that is
      \(\lambda(B)=\nu(B)/\nu(D)\), is a probability measure on it;
\item \(\mu|_D\ll\lambda\) on \((D,\mathcal F_D)\);
\item for every \(A\in\mathcal F\), the set \(A\cap D\) lies in
      \(\mathcal F_D\) and \(\mu|_D(A\cap D)=\mu(A\cap D)\).
\end{enumerate}
Hence the finite measure \(\mu|_D\), the probability measure \(\lambda\)
and the generating sequence \((E_j\cap D)_{j\geq1}\) satisfy every
hypothesis of the finite case of
\cref{thm:p15-countably-generated-radon-nikodym} on the space \(D\).
\end{hbcorollary}

The corollary is applied to each block \(D_i\) of the sigma-finite
assembly with \(\nu(D_i)>0\); blocks with \(\nu(D_i)=0\), including empty
ones, receive the density zero, since then \(\mu(D_i)=0\). The local
densities are measurable for \(\mathcal F_{D_i}\), and part~(4) of
\cref{thm:p15-w02c-p15-003-trace-generation}, with \(S=[0,\infty]\)
carrying the trace on \([0,\infty]\) of the sigma-algebra of
\cref{def:p15-extended-random-variables}, makes the glued density
measurable for \(\mathcal F\).

The finite case delivers each local density as an integrand on the block
\((D,\mathcal F_D)\) against the probability measure \(\lambda\), while
the representation to be proved concerns integrals over \(\Omega\) against
\(\nu\). The next lemma converts the one kind of integral into the other.
For a function \(g:D\to[0,\infty]\), its \emph{zero extension}
\(\bar g:\Omega\to[0,\infty]\) equals \(g\) on \(D\) and \(0\) on
\(\Omega\setminus D\). Integrals on the block are those of
\cref{def:p15-simple-integral} and \cref{def:p15-nonnegative-integral},
formed on the measure space \((D,\mathcal F_D,\lambda)\) or
\((D,\mathcal F_D,\nu|_D)\) in place of \((\Omega,\mathcal F,\mu)\);
for \(B\in\mathcal F_D\), \(\int_Bg\,d\lambda\) is the integral of the
function equal to \(g\) on \(B\) and to \(0\) on \(D\setminus B\).

\begin{hblemma}{thm:p15-w02c-p15-003-block-integrals}{proved=annex:p15-w02c-p15-003-block-integrals}{Integrals on a block and on the whole space}
Let \((\Omega,\mathcal F,\nu)\) be a measure space, let
\(D\in\mathcal F\) be nonempty, and let \(g:D\to[0,\infty]\) be
measurable from \((D,\mathcal F_D)\) into \(\overline{\mathbb R}\).
\begin{enumerate}
\item The zero extension \(\bar g\) is measurable from
      \((\Omega,\mathcal F)\), and
      \(\int_Dg\,d(\nu|_D)=\int_\Omega\bar g\,d\nu\).
\item If \(0<\nu(D)<\infty\) and \(\lambda=\nu|_D/\nu(D)\), then for
      every \(B\in\mathcal F_D\), which is also a member of
      \(\mathcal F\),
      \[
       \int_Bg\,d\lambda=\frac1{\nu(D)}\int_B\bar g\,d\nu .
      \]
\item Under the hypotheses of part~(2), a set \(N\in\mathcal F_D\)
      satisfies \(\lambda(N)=0\) exactly when \(\nu(N)=0\). In
      particular, if \(g<\infty\) holds \(\lambda\)-almost everywhere on
      \(D\), then \(\{\omega\in D:g(\omega)=\infty\}\) is a member of
      \(\mathcal F\) of \(\nu\)-measure zero.
\end{enumerate}
\end{hblemma}

\begin{hbtheorem}{thm:p15-w02c-p15-003-assembly}{proved=annex:p15-w02c-p15-003-assembly}{Assembly of local densities}
Let \(\mu,\nu\) be measures on \((\Omega,\mathcal F)\) with \(\mu\ll\nu\),
and let \((D_i)_{i\in I}\) be a countable family of pairwise disjoint
members of \(\mathcal F\) with union \(\Omega\) and \(\nu(D_i)<\infty\)
for every \(i\). Let \(I^+=\{i\in I:\nu(D_i)>0\}\). For each
\(i\in I^+\), let \(\lambda_i=\nu|_{D_i}/\nu(D_i)\) and let
\(f_i:D_i\to[0,\infty]\) be measurable from \((D_i,\mathcal F_{D_i})\),
finite \(\lambda_i\)-almost everywhere, with
\[
 \int_Bf_i\,d\lambda_i=\mu(B)\qquad(B\in\mathcal F_{D_i}).
\]
Define \(f:\Omega\to[0,\infty]\) by \(f=f_i/\nu(D_i)\) on \(D_i\) for
\(i\in I^+\) and \(f=0\) on \(D_i\) for \(i\in I\setminus I^+\). Then
\(f\) is measurable from \((\Omega,\mathcal F)\), finite
\(\nu\)-almost everywhere, and for every \(A\in\mathcal F\) and every
\(i\in I\),
\[
 \int_{A\cap D_i}f\,d\nu=\mu(A\cap D_i),
 \qquad
 \int_Af\,d\nu=\sum_{i\in I}\mu(A\cap D_i)=\mu(A).
\]
\end{hbtheorem}

When \(\mu(D_i)<\infty\) for every \(i\) and \(\mathcal F\) is generated
by \(E_1,E_2,\ldots\), the corollary \cref{thm:p15-w02c-p15-003-block-data}
makes the finite case of \cref{thm:p15-countably-generated-radon-nikodym}
applicable on every block with \(i\in I^+\), and its conclusion is a
function \(f_i\) with exactly the properties assumed in the theorem. The
theorem then gives the sigma-finite representation and the
finite-almost-everywhere property of the glued density.

\begin{hbexample}{def:p15-w02c-p15-003-trace-example}{A four-point space cut into two blocks}
Let \(\Omega=\{1,2,3,4\}\), \(E_1=\{1,2\}\), \(E_2=\{1,3\}\), and \(E_j=\varnothing\) for \(j\geq3\). The
sets \(E_1\cap E_2=\{1\}\), \(E_1\setminus E_2=\{2\}\),
\(E_2\setminus E_1=\{3\}\) and \(\Omega\setminus(E_1\cup E_2)=\{4\}\)
lie in \(\sigma(E_1,E_2,\ldots)\), so every subset of \(\Omega\) does, being a
finite union of these singletons, and
\(\mathcal F=\sigma(E_1,E_2,\ldots)=2^\Omega\). Take the block
\(D=\{1,2\}\). Its trace is \(\mathcal F_D=2^D\). The restricted
generators are \(E_1\cap D=D\), \(E_2\cap D=\{1\}\) and \(E_j\cap D=\varnothing\) for \(j\geq3\). A sigma-algebra
on \(D\) containing \(\{1\}\) also contains its complement relative to
\(D\), which is \(\{2\}\), so
\(\sigma_D(D,\{1\},\varnothing,\varnothing,\ldots)=\{\varnothing,\{1\},\{2\},D\}=2^D=\mathcal F_D\),
as part~(3) of \cref{thm:p15-w02c-p15-003-trace-generation} asserts.

Let \(\nu(\{k\})=1\) for each \(k\), and let \(\mu(\{1\})=2\),
\(\mu(\{2\})=0\), \(\mu(\{3\})=\mu(\{4\})=1\), each extended to
\(2^\Omega\) by addition. Then \(\mu\ll\nu\), since only the empty set is
\(\nu\)-null, and \(\nu(D)=2\). The probability \(\lambda=\nu|_D/2\)
gives each of \(1\) and \(2\) mass \(1/2\). On \(D\), the partition of
\(D\) into the nonempty Boolean atoms of the restricted generators is
\(\{\{1\},\{2\}\}\), and its ratios give the local density
\(f_D(1)=\mu(\{1\})/\lambda(\{1\})=4\) and \(f_D(2)=0\) relative to
\(\lambda\). The assembly divides by \(\nu(D)=2\), giving \(f(1)=2\) and
\(f(2)=0\) relative to \(\nu\); these are the ratios
\(\mu(\{k\})/\nu(\{k\})\). The other block \(\{3,4\}\) gives
\(f(3)=f(4)=1\) in the same way, with its own restricted generators
\(E_1\cap\{3,4\}=\varnothing\) and \(E_2\cap\{3,4\}=\{3\}\).

The transfer of \cref{thm:p15-w02c-p15-003-block-integrals} can be checked
on \(B=\{1\}\). On the block,
\(\int_Bf_D\,d\lambda=4\cdot\lambda(\{1\})=2\), which is \(\mu(\{1\})\).
The zero extension \(\bar f_D\) takes the values \(4,0,0,0\) at
\(1,2,3,4\), so \(\int_B\bar f_D\,d\nu=4\) and
\(\nu(D)^{-1}\int_B\bar f_D\,d\nu=2\), as part~(2) asserts. For
\(A=\{1,3\}\), the per-block identities of
\cref{thm:p15-w02c-p15-003-assembly} read
\(\int_{A\cap D}f\,d\nu=f(1)=2=\mu(\{1\})\) and
\(\int_{A\cap\{3,4\}}f\,d\nu=f(3)=1=\mu(\{3\})\), and their sum is
\(3=\mu(A)\). Had the local density been glued without the division by
\(\nu(D)=2\), the first block would contribute \(4\), which is
\(\nu(D)\mu(\{1\})\) and not \(\mu(\{1\})\).
\end{hbexample}

\begin{hbwarning}{def:p15-w02c-p15-003-trace-limits}{Complements, nonmeasurable blocks and blocks of infinite mass}
In \cref{def:p15-w02c-p15-003-trace-sigma-field}, \(\sigma_D\) takes
complements relative to \(D\). Generating the same sets inside
\(\Omega\) gives a different collection. In the example of
\cref{def:p15-w02c-p15-003-trace-example}, the sigma-algebra on
\(\Omega\) generated by \(D=\{1,2\}\) and \(\{1\}\) has the atoms
\(\{1\}\), \(\{2\}\) and \(\{3,4\}\); it contains \(\{3,4\}\) and
\(\{2,3,4\}\), which are not subsets of \(D\), so it is not a
sigma-algebra on \(D\) and cannot serve as the sigma-field of a block.

Part~(2) of \cref{thm:p15-w02c-p15-003-trace-generation} and the
corollary need \(D\in\mathcal F\). Let \(\Omega=\{1,2,3\}\),
\(\mathcal F=\{\varnothing,\{1,2\},\{3\},\Omega\}\) and
\(D=\{1,3\}\notin\mathcal F\). Then
\(\mathcal F_D=\{\varnothing,\{1\},\{3\},D\}\), a sigma-algebra on
\(D\) by part~(1), but its member \(\{1\}\) is not in \(\mathcal F\). A
measure on \(\mathcal F\) assigns no value to \(\{1\}\), so no restriction
\(\mu|_D\) is defined. The disjointified blocks of the sigma-finite
assembly are built from members of \(\mathcal F\) by finitely many
intersections and differences, so they lie in \(\mathcal F\) by
\cref{thm:p15-generated-events}.

Parts~(2) and~(3) of \cref{thm:p15-w02c-p15-003-block-integrals} need
\(0<\nu(D)<\infty\). For counting measure \(\nu\) on
\(\Omega=\{1,2,3,\ldots\}\) with \(\mathcal F=2^\Omega\) and
\(D=\Omega\), the mass \(\nu(D)\) is infinite, no positive real
multiple of \(\nu|_D\) has total mass \(1\), and there is no
probability \(\lambda\) to which the finite case could be applied. This
is why the sigma-finite assembly first cuts \(\Omega\) into blocks of
finite \(\nu\)-mass; on a block of mass zero the normalisation is
undefined as well, and the theorem assigns the density zero there instead.
\end{hbwarning}

\emph{Relationships.} The theorem and corollary of this section are the
block step of the sigma-finite assembly in
\cref{thm:p15-countably-generated-radon-nikodym}: part~(3) supplies the
generating sequence on each block, the corollary supplies the finite
measure, the probability measure and the absolute continuity required by
the finite case, and part~(4) gives measurability of the glued density.
\Cref{thm:p15-w02c-p15-003-block-integrals} converts the local
representation on each block into an identity for integrals against
\(\nu\) on \(\Omega\), and \cref{thm:p15-w02c-p15-003-assembly}
sums these identities over the blocks and passes the finiteness of the
local densities from \(\lambda_i\)-null sets to \(\nu\)-null sets.
Parts~(1) and~(2) also describe the trace sigma-algebra used when the
proof of \cref{thm:p15-radon-nikodym} applies its finite construction to
each piece of a sigma-finite partition. The proofs use only
\cref{def:p15-sigma-algebra}, \cref{thm:p15-generated-events},
\cref{def:p15-measure} and \cref{def:p15-measurable-map}; the lemma and
the assembly theorem use in addition \cref{def:p15-extended-random-variables},
the integrals of \cref{def:p15-simple-integral} and
\cref{def:p15-nonnegative-integral}, and \cref{thm:p15-nonnegative-integral-laws}.

\emph{Source.} Durrett's Exercise~4.3.7 asks for the density theorem on a
countably generated sigma-field for sigma-finite measures with
\(\mu\ll\nu\), by applying a martingale convergence theorem to the
partition example \cite[exercise~4.3.7, p.~234]{DurrettPTEAuthor2019}. The
exercise does not discuss the passage from the whole space to blocks of
finite mass; the statements of this section are general measure theory
written for this handbook, not a presentation of the source.
